\section{Methodology}
\label{methodology}

\subsection{Framework Overview}
The framework takes an incomplete natural-language requirement as input and produces a clarified requirement statement with its elicitation record.
It follows a lightweight loop: detect missing information, assign clarification responsibilities to agents, ask knowledge-guided follow-up questions, and let humans validate the refined requirement.
The scope is limited to requirement elicitation.

\subsection{Multi-Agent Requirement Clarification Workflow}
The clarification workflow uses three complementary agents.

\subsubsection{Interviewer Agent}
The interviewer agent manages the elicitation dialogue.
It asks concise follow-up questions, adapts the question order based on stakeholder responses, and avoids unrelated prompts.

\subsubsection{End-User Agent}
The end-user agent represents usage perspectives during clarification.
It helps expose user goals, operational contexts, expected interactions, and practical constraints that may be absent from the original requirement.

\subsubsection{Analyst Agent}
The analyst agent reviews the evolving requirement statement.
It checks whether the current version contains sufficient actor, action, object, condition, constraint, data, and exception information for elicitation.

\subsection{Missing Information Detection}
Missing information detection identifies which requirement elements must be clarified before the requirement can be considered complete for the elicitation task.

\subsubsection{Functional Information}
Functional information covers actors, goals, actions, inputs, outputs, triggers, and expected system responses.

\subsubsection{Constraint Information}
Constraint information covers business rules, quality constraints, timing limits, access restrictions, and usage conditions.

\subsubsection{Role, Data, and Exception Information}
Role, data, and exception information is treated as high priority because these gaps often determine whether stakeholders interpret a requirement consistently.

\subsection{Knowledge-Assisted Clarification}
Knowledge assists the agents in selecting what to ask and how to ask it.

\subsubsection{Requirement Knowledge}
Requirement knowledge includes domain terms, stakeholder roles, requirement slots, constraint categories, and elicitation checklists.

\subsubsection{Clarification Templates}
Clarification templates map missing information types to targeted follow-up questions.
For example, a missing actor triggers a role question, while a missing exception triggers an alternative-flow question.

\subsubsection{Prompt Strategies}
Prompt strategies constrain agents to ask one focused question at a time, justify the missing information type, and avoid introducing requirements not supported by stakeholder input.

\subsection{Human-in-the-Loop Interaction}
Human participants validate clarification questions, provide answers, and approve revised requirement statements.
The workflow preserves human control over requirement meaning while using agents to reduce missed clarification opportunities.

\subsection{Iterative Requirement Refinement}
The workflow repeats until the remaining missing information is resolved or explicitly marked as unknown by the stakeholder.
Each iteration updates the current requirement statement and records the question-answer pair that motivated the refinement.
